%*********************************************************************
\chapter{Further Work}
\label{ch:furtherwork}
%*********************************************************************

\textbf{Heat the plate from below under closed loop, and map the plate.} Run 3 showed a base heater lifts the well to within 0.19\,\textdegree{}C of target with the air unchanged, and run 4 showed an open loop base leaves the chamber loop at 4\,\% duty with nothing left to trim. The base needs its own PI loop with a probe at the plate, which also allows 37\,\textdegree{}C to be held at the sample instead of through the air. A loop sets the average plate temperature but not the 1.81\,\textdegree{}C spread, which is set by where heat enters, so the busbars need moving to opposite long edges or the pane splitting into two zones. Only two wells were measured, so the temperature of every well should then be mapped with the base at setpoint. This would show which wells hold 37\,\textdegree{}C within tolerance and which do not, and if the spread is confined to the outer rows, the edge wells can simply be left empty and the remainder of the plate used. This addresses R2 and R3 together.

\textbf{Supply the humidity actively, through a pumped air circuit.} Only about a twentieth of the air entering the chamber comes through the regulator, so humidifying that stream cannot hold the chamber near saturation however wet the gas arrives. A separate vessel of 5\,\% CO\textsubscript{2} in air, pumped continuously through heated water, sets the flow with a pump instead of the control valve, and adds a sensor and a loop of its own in the way the commercial systems do. It also keeps standing water out of the enclosure, which the BSL-3 work this incubator is intended for does not allow. A pumped circuit is needed for that work regardless of humidity. A chamber sealed for containment no longer exchanges air with the room through its leak, so fresh gas has to be supplied deliberately, and the air leaving the chamber has to be filtered or returned to the circuit instead of escaping into the room. Building the humidified supply as a closed, pumped air circuit therefore solves R3 now and provides the airflow path the sealed BSL-3 version will need.

\textbf{Run the gas loop.} The joint that failed should be avoided by using either a regulator without a filter in its inlet or a standard CO\textsubscript{2} cylinder in place of the SodaStream adapter. Before the first run, the cable gland and gas line should be sealed and the needle valve set to 0.1\,L/min or less, since a faster flow fills the chamber only slightly faster at a large cost in phase margin (Section~\ref{sec:co2-sensitivity}). The leak should then be measured first, by pulsing the valve and fitting a single exponential to the decay that follows, which gives $Q_L = V/\tau$. The controller gains can then be recalculated from Equation~\ref{eq:lambda-gains}, and if the leak turns out high, the control period should be lengthened to limit the ripple.

\textbf{Replace the jumper wiring with a printed circuit board.} Every failure during testing was a wiring connection: two thermistor leads, one of them on the glass, the SCD30, and the I\textsuperscript{2}C connection that ended the long run at 6\,h\,33\,min. No loop tripped, no sensor gave a false reading and the PT1000 reported no faults, so R8 failed because of wiring and not control. A soldered board with secure connectors should be built before another unattended run, with a logic-level MOSFET in place of the TIP31C, which drops 2\,V at 2\,A and wastes 4\,W as heat.

\textbf{Revise the enclosure.} Section~\ref{sec:build-observations} lists four changes the first build showed were needed: screws through the lid frame, since nothing currently holds the lid down; a larger seat for the CO\textsubscript{2} sensor; a second cable gland, since one cannot take four thermistors, the PT1000, two gas sensors and the glass leads; and thicker plate clips with the notch removed and the root filleted, since one snapped under thumb pressure. The 0.8\,mm floor and the 21.9\,mm gap above the plate are both larger than they need to be, so the chamber can be made shorter once the clips are stiffened, which recovers working distance under R7.

\textbf{Test the incubator on a microscope.} Transmission through the ITO lid was never measured. The lid is removable, so it should be compared directly against a plain sheet of the same glass. Focal drift was not measured either, and since none of the reviewed systems report it (Section~\ref{subsec:existing-drift}), a long run on the stage at setpoint would give a result not yet available in that literature. The heated base pane also sits in the light path as built, so a heated metal frame holding the plate should be tested as a way to heat from below with the optical path clear.

\textbf{Recheck the instrumentation, then run a culture.} The thermistor offsets were measured at 23.6\,\textdegree{}C and the PT1000 was calibrated at 0\,\textdegree{}C, so both calibrations should be repeated near the 37\,\textdegree{}C operating point. The humidity sensor was never checked against a reference, and a saturated salt solution in a properly sealed container should be used to provide one. These checks should be written up as a commissioning procedure that the next user can repeat. Once they hold, the incubator should be tested for the purpose it was built for: \textit{M. smegmatis} on an agar pad, imaged continuously for long enough to watch a microcolony form. This is the only test that shows the environment the chamber holds is the one the biology needs.

%   Last page of the body, for the automatic page count on the declaration.
\label{body:end}
